\documentclass[11pt,reqno]{amsart}
\usepackage{iftex}
\ifPDFTeX
  \usepackage[T1]{fontenc}
  \usepackage[utf8]{inputenc}
  \usepackage{lmodern}
\else
  \usepackage[no-math]{fontspec}
\fi
\usepackage{amsmath}
\usepackage{amssymb}
\usepackage{amsfonts}
\usepackage{mathrsfs}
\usepackage{amsthm}
\usepackage{mathtools}
\usepackage{geometry}
\usepackage{enumitem}
\usepackage{longtable}
\usepackage{booktabs}
\usepackage{array}
\usepackage[dvipsnames]{xcolor}
\usepackage[pagebackref=true]{hyperref}
\usepackage{bookmark}
\usepackage{pdfcomment}

% Hover previews for internal references (no layout change).
\makeatletter
\let\CG@origref\ref
\let\CG@origeqref\eqref
\newcommand{\CG@gettip}[1]{\@ifundefined{CGTip@#1}{Reference: #1}{\csname CGTip@#1\endcsname}}
\DeclareRobustCommand{\CGref}[1]{\pdftooltip{\CG@origref{#1}}{\CG@gettip{#1}}}
\DeclareRobustCommand{\CGeqref}[1]{\pdftooltip{\CG@origeqref{#1}}{\CG@gettip{#1}}}
% ---- tooltips for the labels of this insert (copy into the v57 preamble) ----
\expandafter\gdef\csname CGTip@sec:BM-transport\endcsname{Section: The transported product and the identification of its endpoint.}
\expandafter\gdef\csname CGTip@eq:BM-Plambda\endcsname{Equation: P\_lambda(s) = product over k of (1 + s(s-1)/(1/4 + gamma\_k(lambda)\textasciicircum\{\}2)), gamma\_k(lambda) = (1-lambda) gamma\_k + lambda tau\_k.}
\expandafter\gdef\csname CGTip@thm:BM-product\endcsname{Theorem - Transported product: locally uniform convergence on [0,1] x C, zeros 1/2 +- i gamma\_k(lambda), symmetries, P\_lambda(0)=P\_lambda(1)=1, convergence of all derivatives as lambda -> 0.}
\expandafter\gdef\csname CGTip@eq:BM-zeta-lambda\endcsname{Equation: zeta\_lambda = P\_lambda/G, zeta\_1 = zeta\_RVM, zeta\_Z1 = zeta\_0 = P\_0/G.}
\expandafter\gdef\csname CGTip@thm:BM-flow\endcsname{Theorem - Flow: zeta\_lambda -> zeta\_Z1 locally uniformly (spherical metric); d/dlambda log P\_lambda = V\_lambda; exp(-integral\_0\textasciicircum\{\}1 V) = P\_0/P\_1.}
\expandafter\gdef\csname CGTip@eq:BM-generator\endcsname{Equation: V\_lambda(s) = -2s(s-1) sum\_k gamma\_k(lambda) eta\_k / ((1/4+gamma\_k(lambda)\textasciicircum\{\}2)(1/4+gamma\_k(lambda)\textasciicircum\{\}2+s(s-1))).}
\expandafter\gdef\csname CGTip@eq:BM-integrated\endcsname{Equation: zeta\_Z1 = zeta\_RVM P\_0/P\_1 = zeta\_RVM exp(-integral\_0\textasciicircum\{\}1 V\_u du).}
\expandafter\gdef\csname CGTip@eq:BM-xi-product\endcsname{Equation: 2 xi(s) = product over rho with Im rho > 0 of (1 + s(s-1)/(rho(1-rho))).}
\expandafter\gdef\csname CGTip@prop:BM-normalization\endcsname{Proposition - Normalization of the endpoint: P\_0 = 2 xi iff every zero has beta = 1/2; equivalently zeta\_Z1 = zeta iff RH.}
\expandafter\gdef\csname CGTip@thm:BM-identification\endcsname{Theorem - Identification: M\_Z1 = M iff 1/zeta\_Z1 has Dirichlet coefficients mu(n) iff zeta\_Z1 = zeta iff RH.}
\expandafter\gdef\csname CGTip@prop:BM-no-Dirichlet\endcsname{Proposition - Neither zeta\_RVM nor 1/zeta\_RVM is an absolutely convergent Dirichlet series in any half-plane: 1/zeta\_RVM(sigma+it) -> L as |t| -> infinity and Carlson's mean value theorem.}
\expandafter\gdef\csname CGTip@rem:BM-status\endcsname{Remark - Status: the transport and the equivalence are proved; M\_Z1 = M is not; the endpoint depends only on the ordinates.}
\makeatother

\hypersetup{
  colorlinks=false,
  pdfborder={0 0 1},
  linkbordercolor={1 0 0},
  citebordercolor={0 0.3 1},
  urlbordercolor={0 0.6 0.6},
  filebordercolor={0 0.6 0.6},
  linktoc=all,
  bookmarksnumbered=true,
  bookmarksopen=true,
  bookmarksopenlevel=1,
  pdfstartview=FitH,
  pdftitle={Bounded M_RVM: transport of ordinates and identification of the endpoint},
  pdfauthor={V. Casa\~na-Giner},
  pdfsubject={Continuous transport of the midpoint RVM product to the genuine ordinates; identification of the endpoint with zeta and its status},
  pdfkeywords={Riemann zeta function, Riemann-von Mangoldt clock, Hadamard product, Mertens function, transport of ordinates}
}
% Back-references in the bibliography.
\renewcommand*{\backref}[1]{}
\renewcommand*{\backrefalt}[4]{%
  \ifcase #1 %
    \relax
  \or
    \ (Cited on p.~#2.)%
  \else
    \ (Cited on pp.~#2.)%
  \fi}

% In \cite[note]{key} the note (e.g. a section number) is linked too.
\makeatletter
\let\CG@origcite\cite
\renewcommand*{\cite}[2][]{\def\CG@citekey{#2}%
  \ifx\relax#1\relax\CG@origcite{#2}\else\CG@origcite[#1]{#2}\fi}
\def\@cite#1#2{{\@citestyle[\citeform{#1}\if@tempswa, \hyper@link{cite}{cite.\CG@citekey}{#2}\fi]}}
\makeatother

\geometry{margin=1in}
\numberwithin{equation}{section}
\setcounter{tocdepth}{3}
\setlength{\emergencystretch}{1.5em}

\theoremstyle{plain}
\newtheorem{theorem}{Theorem}[section]
\newtheorem{proposition}[theorem]{Proposition}
\newtheorem{lemma}[theorem]{Lemma}
\newtheorem{corollary}[theorem]{Corollary}
\theoremstyle{definition}
\newtheorem{definition}[theorem]{Definition}
\newtheorem{hypothesis}[theorem]{Hypothesis}
\newtheorem{heuristic}[theorem]{Heuristic}
\theoremstyle{remark}
\newtheorem{remark}[theorem]{Remark}
% As in v56:
\theoremstyle{plain}
\newtheorem{auditproposition}{Audit proposition}
\renewcommand{\theauditproposition}{A\arabic{auditproposition}}

\begin{document}

\title[Bounded $M_{\rm RVM}$]{Bounded $M_{\rm RVM}$: transport of ordinates\\
and identification of the endpoint}
\author{V. Casa\~na-Giner}
\date{4 October 2026}
\maketitle

\noindent\textbf{Scope.} This note reviews and corrects the three-step
argument that transports the critical RVM product to the genuine
ordinates.  Steps~1 and~2 are correct and are stated with proofs; the
identification theorem of Step~3 is correct as an equivalence; the
verification proposed in Step~3 through Dirichlet coefficients of
$1/\zeta_{\rm RVM}$ is not available, because those coefficients do not
exist (Proposition~\CGref{prop:BM-no-Dirichlet}).  The notation and the
transition theorem (TR) are those of the insert
\texttt{TE\_tracking\_v57}.

\tableofcontents

\section{The transported product and the identification of its endpoint}
\label{sec:BM-transport}

Let $\gamma_k$ be the $k$-th genuine ordinate, counted with multiplicity
and without assuming $\beta_k=\tfrac12$; let $\tau_k=\gamma_k^{\rm RVM}$,
defined by $\mathcal N_{\rm RVM}(\tau_k)=k-\tfrac12$, and
$\eta_k=\tau_k-\gamma_k$.  Put $q(s)=s(s-1)$ and
$G(s)=s(s-1)\Gamma(s/2)\pi^{-s/2}$.

\subsection{Step 1: the transported product}

For $0\le\lambda\le1$ let $\gamma_k(\lambda)=(1-\lambda)\gamma_k+\lambda\tau_k$ and
\begin{equation}\label{eq:BM-Plambda}
P_\lambda(s)=\prod_{k\ge1}\Big(1+\frac{s(s-1)}{\tfrac14+\gamma_k(\lambda)^2}\Big).
\end{equation}

\begin{theorem}[Transported product]\label{thm:BM-product}
The product \eqref{eq:BM-Plambda} converges locally uniformly on
$[0,1]\times\mathbb C$.  Each $P_\lambda$ is entire, depends continuously on
$\lambda$, has zeros exactly at $\tfrac12\pm i\gamma_k(\lambda)$ (with
multiplicity), satisfies $P_\lambda(1-s)=P_\lambda(s)$,
$P_\lambda(\bar s)=\overline{P_\lambda(s)}$ and $P_\lambda(0)=P_\lambda(1)=1$,
and $P_\lambda^{(j)}\to P_0^{(j)}$ locally uniformly as $\lambda\to0^+$ for
every $j\ge0$.  Moreover $P_1=P^{\rm crit}_{\rm RVM}$.
\end{theorem}

\begin{proof}
From $N(T)=\mathcal N_{\rm RVM}(T)+S(T)+O(1/T)$ and $S(T)=O(\log T)$ one
gets $\mathcal N_{\rm RVM}(\tau_k)-\mathcal N_{\rm RVM}(\gamma_k)=O(\log\gamma_k)$;
since $\mathcal N'_{\rm RVM}(t)=\frac1{2\pi}\log\frac t{2\pi}$, the mean value
theorem gives $\eta_k=O(1)$, and $\gamma_k\asymp\tau_k\asymp k/\log k$.  On a
compact set $K$, $|q|\le C_K$ and $\tfrac14+\gamma_k(\lambda)^2\gg\gamma_k^2$
uniformly in $\lambda$ for $k$ large, so
$\sup_{K\times[0,1]}|q/(\tfrac14+\gamma_k(\lambda)^2)|\ll_K(\log k)^2k^{-2}$,
which is summable: the product converges normally.  The zero statement
follows from
$\tfrac14+\gamma_k(\lambda)^2+s(s-1)=(s-\tfrac12)^2+\gamma_k(\lambda)^2$, the
symmetries from the dependence on $s(s-1)$ and the reality of
$\gamma_k(\lambda)$, and $P_\lambda(0)=P_\lambda(1)=1$ from $q(0)=q(1)=0$.
With $A_{k,\lambda}=\tfrac14+\gamma_k(\lambda)^2$,
$|q/A_{k,\lambda}-q/A_{k,0}|\ll_K\lambda|\eta_k|\gamma_k^{-3}+\lambda^2\eta_k^2\gamma_k^{-4}$,
which is summable; hence $P_\lambda\to P_0$ locally uniformly, and Cauchy's
formula gives the derivatives.
\end{proof}

\subsection{Step 2: the flow of $\zeta_{\rm RVM}$}

Define
\begin{equation}\label{eq:BM-zeta-lambda}
\zeta_\lambda=\frac{P_\lambda}{G},\qquad \zeta_1=\zeta_{\rm RVM},\qquad
\zeta_{Z1}:=\zeta_0=\frac{P_0}{G},\qquad F_\lambda=\frac1{\zeta_\lambda}.
\end{equation}

\begin{theorem}[Flow]\label{thm:BM-flow}
As $\lambda\to0^+$, $\zeta_\lambda\to\zeta_{Z1}$ locally uniformly in the
spherical metric, and $F_\lambda\to F_0$ uniformly on compact sets free of
zeros of $P_0$.  Off the moving divisor,
\begin{equation}\label{eq:BM-generator}
\partial_\lambda\log P_\lambda(s)=\mathcal V_\lambda(s):=
-2s(s-1)\sum_{k\ge1}\frac{\gamma_k(\lambda)\eta_k}
{\big(\tfrac14+\gamma_k(\lambda)^2\big)\big(\tfrac14+\gamma_k(\lambda)^2+s(s-1)\big)},
\end{equation}
the series converging locally uniformly, and
\begin{equation}\label{eq:BM-integrated}
\zeta_{Z1}=\zeta_{\rm RVM}\,\frac{P_0}{P_1}
=\zeta_{\rm RVM}\exp\Big(-\int_0^1\mathcal V_u\,du\Big).
\end{equation}
If a zero $\tfrac12+i\gamma_k$ of $P_0$ is simple, the residues
$G(\rho_{k,\lambda})/(\rho_{k,\lambda}P'_\lambda(\rho_{k,\lambda}))$ of
$F_\lambda(s)/s$ at $\rho_{k,\lambda}=\tfrac12+i\gamma_k(\lambda)$ converge to
the corresponding residue of $F_0(s)/s$.
\end{theorem}

\begin{proof}
Uniform convergence of $P_\lambda$ and division by the fixed function $G$
give the first claim; on a compact set with $\min|P_0|=m>0$ one has
$|P_\lambda|\ge m/2$ for small $\lambda$ and
$|F_\lambda-F_0|\le2\sup|G|\,m^{-2}|P_\lambda-P_0|$.  Termwise,
$\partial_\lambda\operatorname{Log}(1+q/A_{k,\lambda})=-2q\gamma_k(\lambda)\eta_k/(A_{k,\lambda}(A_{k,\lambda}+q))
=O_K(|\eta_k|\gamma_k^{-3})$, summable.  For \eqref{eq:BM-integrated},
integrating each mode with $u=\gamma_k(\lambda)$, $du=\eta_k\,d\lambda$, gives
$\frac1{2q}\operatorname{Log}\frac{u^2+1/4}{u^2+1/4+q}\big|_{\gamma_k}^{\tau_k}$
modulo $2\pi i$ (the principal branch is valid only while the argument
does not cross $(-\infty,0]$); after exponentiation the factor
$\prod_k(1+q/A_{k,0})/(1+q/A_{k,1})=P_0/P_1$ is exact.  The residues
converge by Theorem~\CGref{thm:BM-product}.
\end{proof}

Identity \eqref{eq:BM-integrated} is the trivial factorization
$P_0/G=(P_1/G)(P_0/P_1)$ written as a flow.  The endpoint is determined
directly by the genuine ordinates placed on the critical line:
\[
\zeta_{Z1}(s)=\frac1{G(s)}\prod_{k\ge1}\Big(1+\frac{s(s-1)}{\tfrac14+\gamma_k^2}\Big).
\]

\subsection{Step 3: identification of the endpoint}

The Hadamard factorization of
$\xi(s)=\tfrac12s(s-1)\pi^{-s/2}\Gamma(s/2)\zeta(s)$, pairing $\rho$ with
$1-\rho$ and using $2\xi(0)=1$, gives
\begin{equation}\label{eq:BM-xi-product}
2\xi(s)=\prod_{\rho:\ \Im\rho>0}\Big(1+\frac{s(s-1)}{\rho(1-\rho)}\Big),
\end{equation}
the product taken in order of increasing $\Im\rho$.

\begin{proposition}[Normalization of the endpoint]\label{prop:BM-normalization}
$P_0=2\xi$ if and only if every nontrivial zero has real part $\tfrac12$.
Equivalently, $\zeta_{Z1}=\zeta$ if and only if RH holds.
\end{proposition}

\begin{proof}
For $\rho=\beta+i\gamma$,
$\rho(1-\rho)=\beta(1-\beta)+\gamma^2+i\gamma(1-2\beta)$, which equals
$\tfrac14+\gamma^2$ if and only if $\beta=\tfrac12$.  If RH holds, the
factors of \eqref{eq:BM-xi-product} and of $P_0$ coincide and $P_0=2\xi$.
Conversely, if $P_0=2\xi$ the two functions have the same zeros, and the
zeros of $P_0$ lie on $\Re s=\tfrac12$.  Since $\zeta=2\xi/G$ and
$\zeta_{Z1}=P_0/G$, the second statement follows.  The constant is fixed by
$P_0(0)=1=2\xi(0)$.
\end{proof}

For $c>1$ put
$M_{Z1}(x)=\lim_{U\to\infty}\frac1{2\pi i}\int_{c-iU}^{c+iU}F_0(s)\frac{x^s}{s}\,ds$
whenever the limit exists.

\begin{theorem}[Identification]\label{thm:BM-identification}
The following are equivalent:
\begin{enumerate}[label=\textup{(\roman*)},leftmargin=2.2em]
\item $M_{Z1}$ exists and $M_{Z1}(n)=M(n)$ for every $n\ge1$;
\item $1/\zeta_{Z1}$ is, in $\Re s>1$, the Dirichlet series
$\sum_n\mu(n)n^{-s}$;
\item $\zeta_{Z1}=\zeta$;
\item RH.
\end{enumerate}
\end{theorem}

\begin{proof}
(i)$\Rightarrow$(ii): the normalized Mellin transform
$s\int_1^\infty M(x)x^{-s-1}dx=1/\zeta(s)$ coincides with $1/\zeta_{Z1}$ for
$\Re s>c$.  (ii)$\Rightarrow$(iii): identity of meromorphic functions.
(iii)$\Leftrightarrow$(iv): Proposition~\CGref{prop:BM-normalization}.
(iii)$\Rightarrow$(i): Perron's formula for $1/\zeta$.
\end{proof}

\begin{proposition}[No Dirichlet expansion for $\zeta_{\rm RVM}$]\label{prop:BM-no-Dirichlet}
Neither $\zeta_{\rm RVM}$ nor $1/\zeta_{\rm RVM}$ is an absolutely convergent
Dirichlet series in any half-plane $\Re s>\sigma_0$.
\end{proposition}

\begin{proof}
In the proof of the transition theorem (TR) of \texttt{TE\_tracking\_v57},
the restriction $\sigma'<\tfrac12$ ($s=\tfrac12+\sigma'+it$) is used only in
the passage from $w$ to $\chi$; the form
$\zeta_{\rm RVM}(s)=Z_{\rm RVM}(1+w)(1+O(1/t))$, with
$w=e^{-2\pi i\mathcal N_{\rm RVM}(t-i\sigma')}$, holds uniformly on every
compact range of $\sigma'>0$ (in the bound on the rays it suffices that
$t^2\ge2(\tau_0^2+\sigma'^2)$).  For fixed $\sigma'>0$,
$\Im\mathcal N_{\rm RVM}(t-i\sigma')=-\frac{\sigma'}{2\pi}\log\frac t{2\pi}+O(t^{-2})\to-\infty$,
so $w\to0$ and $1/\zeta_{\rm RVM}(\sigma+it)\to L=1/Z_{\rm RVM}$ as
$|t|\to\infty$, for every fixed $\sigma>\tfrac12$.  If
$1/\zeta_{\rm RVM}-L=\sum_nc(n)n^{-s}$ converged absolutely on
$\Re s=\sigma>\sigma_0$, the mean value theorem for absolutely convergent
Dirichlet series would give
\[
\sum_{n\ge1}|c(n)|^2n^{-2\sigma}=\lim_{T\to\infty}\frac1{2T}\int_{-T}^{T}
\Big|\frac1{\zeta_{\rm RVM}(\sigma+it)}-L\Big|^2dt=0,
\]
so $1/\zeta_{\rm RVM}\equiv L$, contradicting $1/\zeta_{\rm RVM}(1)=0$.  The
same argument with $\zeta_{\rm RVM}\to Z_{\rm RVM}$ excludes
$\zeta_{\rm RVM}$.
\end{proof}

Consequently the coefficients $a_{\rm RVM}(n)$ with
$1/\zeta_{\rm RVM}=\sum a_{\rm RVM}(n)n^{-s}$ do not exist, and under RH
neither do coefficients $\kappa(n)$ of $\exp\int_0^1\mathcal V=\zeta_{\rm RVM}/\zeta$,
since otherwise $\zeta_{\rm RVM}=\zeta\cdot\sum\kappa(n)n^{-s}$ would be a
Dirichlet series.  The proposed check
$(a_{\rm RVM}*\kappa)(n)=\mu(n)$ is therefore not a verifiable condition.

\begin{remark}[Status]\label{rem:BM-status}
Proved: the transport
$\zeta_{\rm RVM}=\zeta_1\to\zeta_\lambda\to\zeta_{Z1}=P_0/G$, locally
uniformly and with divisors and residues transported, and the equivalence
$M_{Z1}=M\iff\zeta_{Z1}=\zeta\iff P_0=2\xi\iff\mathrm{RH}$.  Not proved:
$M_{Z1}=M$.  The endpoint depends only on the ordinates $\gamma_k$, while
the content of RH lies in the abscissae $\beta_k$, which the transport does
not move: $\zeta_{Z1}=\zeta$ exactly when every factor
$1+s(s-1)/(\rho(1-\rho))$ of \eqref{eq:BM-xi-product} equals the
corresponding factor $1+s(s-1)/(\tfrac14+\gamma^2)$ of $P_0$.
\end{remark}

\end{document}
